\section{From Green--Kubo Response to Response Geometry}
\label{sec:kubo_response_geometry}
The preceding sections emphasized thermodynamic response as
the origin of metric and curvature structures in equilibrium and
near-critical systems. We now connect this viewpoint to the
standard language of response theory. The purpose is not to
replace Green--Kubo or nonlinear response theory, but to
reinterpret them geometrically. Green functions, susceptibilities,
response tensors, and nonlinear response functions remain the
central physical objects. Response geometry emerges when one
asks how these response objects themselves vary over a manifold
of external controls.

\subsection{Linear Response and Retarded Green Functions}
Consider a system perturbed by a weak external field \(f(t)\)
coupled to an operator \(B\),
\begin{align}
H(t)=H_0-f(t)B .
\end{align}
To first order in the perturbation, the induced change in the
expectation value of an observable \(A\) is
\begin{align}
\delta\langle A(t)\rangle
=
\int_{-\infty}^{\infty}dt'\,
\chi_{AB}(t-t')f(t'),
\end{align}
where causality requires \(\chi_{AB}(t-t')=0\) for
\(t<t'\). In quantum mechanics the response function is the
retarded Green function
\begin{align}
\chi_{AB}(t)
=
-\frac{i}{\hbar}\Theta(t)
\langle [A(t),B(0)]\rangle_0 .
\label{eq:kubo_response}
\end{align}
Equation~\eqref{eq:kubo_response} is the Kubo formula. It
expresses a measurable susceptibility in terms of an equilibrium
correlation function and therefore provides the microscopic
origin of linear response.
In frequency space,
\begin{align}
\delta\langle A(\omega)\rangle
=
\chi_{AB}(\omega)f(\omega),
\end{align}
so that the susceptibility is the linear map taking a perturbing
field into an observable response. From the perspective of this
review, this is the first geometric step: response is not simply a
number, but a bilinear object relating a direction in perturbation
space to a direction in observable space.

\subsection{From Green--Kubo Response to Geometric Response Theory}
\label{sec:GreenKuboGeometry}
Linear response theory provides the natural starting point for a
geometric description of thermodynamic and quantum response.
Traditionally, the Green--Kubo formalism relates transport
coefficients and susceptibilities to equilibrium correlation
functions, while the fluctuation--dissipation theorem connects
these response functions to spontaneous fluctuations in
equilibrium. The geometric framework developed in this review
extends this picture by showing that the response tensor naturally
decomposes into reciprocal and nonreciprocal sectors, giving rise
to a metric geometry associated with fluctuations and a curvature
geometry associated with cyclic transport.
Consider an open quantum system governed by the parameter-dependent
Liouvillian
%
\begin{equation}
\dot{\rho}
=
\mathcal L(\bm{\lambda})\rho,
\end{equation}
%
with stationary state satisfying
%
\begin{equation}
\mathcal L(\bm{\lambda})
\rho_{\rm ss}
=
0.
\end{equation}
%
For a family of control parameters
$\bm{\lambda}=\{\lambda^\mu\}$,
we define the generalized forces
%
\begin{equation}
O_\mu
=
\frac{\partial H}
{\partial\lambda^\mu},
\end{equation}
%
and the corresponding response one-form
%
\begin{equation}
A_\mu
=
\mathrm{Tr}
\left(
\rho_{\rm ss}
O_\mu
\right)
=
\langle O_\mu\rangle_{\rm ss}.
\end{equation}
%
The exterior derivative of this one-form defines the response
curvature,
%
\begin{equation}
F_{\mu\nu}
=
\partial_\mu A_\nu
-
\partial_\nu A_\mu,
\end{equation}
%
whose surface integral gives the geometric work accumulated over a
closed quasistatic control cycle,
%
\begin{equation}
W
=
\oint A
=
\iint_\Sigma
F.
\end{equation}
To connect this geometric construction with response theory,
differentiate the stationary-state condition,
%
\begin{equation}
\mathcal L
\rho_{\rm ss}
=
0,
\end{equation}
%
to obtain
%
\begin{equation}
\mathcal L
\,
\partial_\mu\rho_{\rm ss}
=
-
(\partial_\mu\mathcal L)
\rho_{\rm ss}.
\end{equation}
%
Introducing the Drazin inverse of the Liouvillian,
%
\begin{equation}
\mathcal L^+,
\end{equation}
%
gives
%
\begin{equation}
\partial_\mu\rho_{\rm ss}
=
-
\mathcal L^+
(\partial_\mu\mathcal L)
\rho_{\rm ss}.
\end{equation}
For Hamiltonian control parameters,
%
\begin{equation}
(\partial_\mu\mathcal L)\rho
=
-i
[O_\mu,\rho],
\end{equation}
%
and using the integral representation
%
\begin{equation}
\mathcal L^+
=
-
\int_0^\infty
dt\,
e^{\mathcal Lt},
\end{equation}
%
one finds
%
\begin{equation}
\partial_\mu\rho_{\rm ss}
=
-i
\int_0^\infty
dt\,
e^{\mathcal Lt}
[O_\mu,\rho_{\rm ss}],
\end{equation}
%
so that
%
\begin{equation}
\partial_\mu A_\nu
=
-i
\int_0^\infty
dt\,
\mathrm{Tr}
\left[
O_\nu
e^{\mathcal Lt}
[O_\mu,\rho_{\rm ss}]
\right].
\end{equation}
The quantum regression theorem identifies the Liouvillian
propagator appearing above with experimentally accessible
two-time correlation functions,
%
\begin{equation}
\mathrm{Tr}
\left[
O_\nu
e^{\mathcal Lt}
(O_\mu\rho_{\rm ss})
\right]
=
\langle
O_\nu(t)
O_\mu(0)
\rangle_{\rm ss},
\end{equation}
%
leading to
%
\begin{equation}
\partial_\mu A_\nu
=
-i
\int_0^\infty
dt\,
\langle
[O_\nu(t),O_\mu(0)]
\rangle_{\rm ss}.
\end{equation}
%
The response curvature therefore becomes
%
\begin{equation}
\boxed{
F_{\mu\nu}
=
-i
\int_0^\infty
dt
\left[
\langle
[O_\nu(t),O_\mu(0)]
\rangle
-
\langle
[O_\mu(t),O_\nu(0)]
\rangle
\right].
}
\label{eq:CurvatureCorrelation}
\end{equation}
Equation~(\ref{eq:CurvatureCorrelation}) establishes that the
geometric curvature is the integrated antisymmetric correlation
function of the generalized forces.
This connection becomes particularly transparent upon introducing
the retarded susceptibility,
%
\begin{equation}
\chi_{\mu\nu}^{R}(t)
=
-i
\theta(t)
\langle
[O_\mu(t),O_\nu(0)]
\rangle,
\end{equation}
%
whose zero-frequency limit is
%
\begin{equation}
\chi_{\mu\nu}^{R}(0)
=
-i
\int_0^\infty
dt\,
\langle
[O_\mu(t),O_\nu(0)]
\rangle.
\end{equation}
%
The response curvature may therefore be written as
%
\begin{equation}
\boxed{
F_{\mu\nu}
=
\chi_{\nu\mu}^{R}(0)
-
\chi_{\mu\nu}^{R}(0),
}
\label{eq:CurvatureKubo}
\end{equation}
%
showing that the curvature is precisely the antisymmetric
zero-frequency response tensor.
Equations~(\ref{eq:CurvatureCorrelation}) and
(\ref{eq:CurvatureKubo}) suggest a natural geometric
decomposition of response theory. The reciprocal sector of the
response tensor is generated by the symmetric correlation
functions,
%
\begin{equation}
G_{\mu\nu}
\sim
\int_0^\infty
dt\,
\left\langle
\{
\delta O_\mu(t),
\delta O_\nu(0)
\}
\right\rangle,
\end{equation}
%
while the nonreciprocal sector is generated by the corresponding
commutator correlations,
%
\begin{equation}
F_{\mu\nu}
\sim
\int_0^\infty
dt\,
\left\langle
[
\delta O_\mu(t),
\delta O_\nu(0)
]
\right\rangle.
\end{equation}
From this perspective, the Green--Kubo formalism naturally
decomposes into two complementary geometric sectors. The symmetric
sector determines a response metric that characterizes
distinguishability, fluctuations, and susceptibility, while the
antisymmetric sector determines a response curvature that
characterizes circulation, geometric work, and nonreciprocal
transport. Response geometry therefore extends the conventional
fluctuation--dissipation theorem by recognizing that reciprocal
and nonreciprocal response are complementary geometric
manifestations of the same underlying correlation functions.


\subsection{Reciprocity, Onsager Symmetry, and Geometric Response}
\label{sec:ReciprocityGeometry}
The preceding discussion shows that response functions admit a
natural geometric interpretation. This perspective becomes
particularly transparent when viewed through the lens of Onsager
reciprocity. In systems satisfying microscopic reversibility, the
static susceptibility is reciprocal,
%
\begin{equation}
\chi_{\mu\nu}
=
\chi_{\nu\mu},
\end{equation}
%
reflecting the underlying symmetry imposed by detailed balance.
The response tensor is therefore entirely symmetric and may be
identified with a Riemannian metric on the control manifold. The
resulting geometry is one of distinguishability, fluctuations,
and susceptibility, and corresponds to the familiar equilibrium
geometries introduced by Weinhold and Ruppeiner.

Away from equilibrium, however, microscopic reversibility need
not hold. External driving, coherent environmental coupling,
nonequilibrium steady states, or measurement backaction may all
break reciprocity, producing an antisymmetric sector of the
response tensor. Rather than viewing this contribution as a small
correction to the reciprocal response, it is more natural to
regard it as defining an independent geometric structure.

Accordingly, we decompose the static response tensor into its
symmetric and antisymmetric components,
%
\begin{equation}
\chi_{\mu\nu}
=
G_{\mu\nu}
+
F_{\mu\nu},
\end{equation}
%
where
%
\begin{align}
G_{\mu\nu}
&=
\frac12
\left(
\chi_{\mu\nu}
+
\chi_{\nu\mu}
\right),
\\
F_{\mu\nu}
&=
\frac12
\left(
\chi_{\mu\nu}
-
\chi_{\nu\mu}
\right).
\end{align}
%
The tensor $G_{\mu\nu}$ defines the reciprocal response metric,
while $F_{\mu\nu}$ defines the antisymmetric response curvature.
These objects are not introduced independently; they are simply
the two irreducible components of the response tensor itself.
This decomposition admits an equally natural interpretation in
terms of correlation functions. As shown in the previous section,
the symmetric sector is generated by the anticommutator
correlations,
%
\begin{equation}
G_{\mu\nu}
\sim
\int_0^\infty
dt\,
\left\langle
\{
\delta O_\mu(t),
\delta O_\nu(0)
\}
\right\rangle,
\end{equation}
%
while the antisymmetric sector is generated by the corresponding
commutator correlations,
%
\begin{equation}
F_{\mu\nu}
\sim
\int_0^\infty
dt\,
\left\langle
[
\delta O_\mu(t),
\delta O_\nu(0)
]
\right\rangle.
\end{equation}
%
The Green--Kubo formalism therefore naturally separates into two
complementary geometric sectors: one associated with equilibrium
fluctuations and reciprocal response, and the other associated
with circulation, geometric work, and nonreciprocal transport.

This viewpoint extends the traditional fluctuation--dissipation
theorem. Rather than identifying response solely with
susceptibility, it recognizes that response possesses both
reciprocal and nonreciprocal components. The reciprocal sector
determines distances on the control manifold through the response
metric, while the antisymmetric sector determines the circulation
of the response one-form through the curvature. Metric geometry
and curvature geometry therefore emerge as complementary
manifestations of the same underlying response tensor.

From this perspective, the appearance of geometric curvature is
not an abstract mathematical construction but the natural
geometric signature of broken reciprocity. Equilibrium systems,
for which Onsager symmetry is preserved, possess purely metric
response geometries. Nonequilibrium systems acquire an additional
antisymmetric sector whose integral over closed loops generates
geometric work. The passage from equilibrium to nonequilibrium
response may therefore be viewed as a transition from purely
Riemannian geometry to a richer differential geometry in which
metric and curvature coexist as complementary descriptors of
physical response.

\subsection{Reciprocity, Fluctuations, and the Symmetric Sector}

In equilibrium, time-reversal symmetry and microscopic
reversibility impose strong constraints on the response matrix.
Onsager reciprocity relates cross-responses associated with
thermodynamically conjugate fluxes and forces, while the
fluctuation--dissipation theorem expresses dissipative response
in terms of spontaneous equilibrium fluctuations. These results
explain why equilibrium response matrices are commonly
symmetric and why their symmetric sector is naturally connected
to covariance, susceptibility, dissipation, and thermodynamic
stability.
This is the response-theory origin of the metric sector discussed
earlier. The same object that appears experimentally as a
susceptibility also appears statistically as a fluctuation and
geometrically as a metric. In schematic form,
\begin{align}
\textrm{fluctuations}
\quad\Longleftrightarrow\quad
\textrm{linear response}
\quad\Longleftrightarrow\quad
\textrm{metric geometry}.
\end{align}
This equivalence underlies the Weinhold and Ruppeiner
constructions, Fisher information geometry, thermodynamic
length, and quantum Fisher metrics. In each case, the metric
describes the local distinguishability or dissipative cost
associated with nearby states or nearby control protocols.
The Green--Kubo viewpoint also clarifies why equilibrium
thermodynamic geometry is integrable. If the generalized forces
are derivatives of a thermodynamic potential, the corresponding
response matrix is a Hessian. Its mixed derivatives commute, and
the antisymmetric sector vanishes. Equilibrium response is
therefore locally reciprocal because it derives from an underlying
state function.
\subsection{Beyond Equilibrium: Antisymmetric Response}
Away from equilibrium, the response matrix need not be purely
symmetric. Persistent currents, broken detailed balance,
dissipation, coherent driving, and environmental basis selection
can all generate response components that are odd under
exchange of control directions. Writing a local response tensor as
\begin{align}
\chi_{\mu\nu}
=
\frac{\partial A_\mu}{\partial\lambda^\nu},
\end{align}
we may decompose it into symmetric and antisymmetric parts,
\begin{align}
G_{\mu\nu}
=
\frac12
\left(
\chi_{\mu\nu}+\chi_{\nu\mu}
\right),
\qquad
F_{\mu\nu}
=
\frac12
\left(
\chi_{\mu\nu}-\chi_{\nu\mu}
\right).
\label{eq:GF_decomp_kubo}
\end{align}
The symmetric part \(G_{\mu\nu}\) is the reciprocal response
metric. It generalizes the fluctuation and friction tensors of
equilibrium and near-equilibrium thermodynamics. The
antisymmetric part \(F_{\mu\nu}\) is the nonreciprocal response
curvature. It measures the failure of response to be integrable
and governs the work or transport accumulated around closed
cycles in control space.
Thus the Green--Kubo framework naturally leads to the
decomposition central to response geometry:
\begin{align}
\boxed{
\textrm{response tensor}
\quad
\chi_{\mu\nu}
=
G_{\mu\nu}
+
F_{\mu\nu}
}
\end{align}
with the symmetric sector describing fluctuations,
distinguishability, and reciprocal susceptibility, and the
antisymmetric sector describing circulation, nonreciprocity, and
holonomy.
This distinction is especially important for nonequilibrium
steady states. Generalized fluctuation--dissipation and
Green--Kubo relations show that response in a nonequilibrium
steady state may still be represented through correlation
functions, but the correlations are no longer those of a detailed
balance equilibrium ensemble. Additional terms associated with
steady currents, entropy production, or local probability currents
enter the response. In this sense, nonequilibrium response theory
already contains the ingredients of curvature: the response
depends on transport structure rather than on fluctuations alone.

\subsection{Response of the Response}
The geometric viewpoint requires one further step. Conventional
response theory asks how an observable changes under an
external perturbation,
\begin{align}
f
\longmapsto
\delta\langle A\rangle .
\end{align}
Response geometry asks how this response changes as the
system itself is moved through a control manifold,
\begin{align}
\lambda
\longmapsto
\chi_{\mu\nu}(\lambda).
\end{align}
Geometry therefore arises from the response of the response.
The hierarchy is
\begin{align}
\boxed{
\mathcal L
\;\longrightarrow\;
G^R(t)
\;\longrightarrow\;
\chi_{\mu\nu}
\;\longrightarrow\;
(G_{\mu\nu},F_{\mu\nu})
}
\end{align}
where the Liouvillian or Hamiltonian generates the dynamics,
the retarded Green function encodes causal propagation, the
susceptibility defines the local response tensor, and the
symmetric and antisymmetric components define the metric and
curvature sectors of response geometry.
This hierarchy is the conceptual bridge between Kubo theory and
the geometric formulation developed in the remainder of the
review. The response tensor is not merely a collection of
coefficients. It is a local geometric object on the control
manifold. Its symmetric part defines a metric-like structure, its
antisymmetric part defines a curvature-like structure, and its
variation over parameter space defines the response landscape.
\subsection{Nonlinear Response and Liouville Space}
The same logic extends naturally to nonlinear spectroscopy. In
optical response theory the macroscopic polarization is expanded
in powers of the applied electric field,
\begin{align}
P(t)
=
P^{(1)}(t)
+
P^{(2)}(t)
+
P^{(3)}(t)
+\cdots .
\end{align}
For third-order spectroscopy,
\begin{align}
P^{(3)}(t)
=
\int_0^\infty dt_3
\int_0^\infty dt_2
\int_0^\infty dt_1\,
R^{(3)}(t_3,t_2,t_1)
E(t-t_3)
E(t-t_3-t_2)
E(t-t_3-t_2-t_1),
\end{align}
where \(R^{(3)}\) is a multi-time response function. In
Liouville space this response is constructed from a sequence of
field interactions and propagations under the system
Liouvillian. The familiar double-sided Feynman diagrams of
nonlinear spectroscopy are therefore a pathway decomposition of
the same causal response structure.
This observation prepares the transition to observational
geometry. In ordinary nonlinear spectroscopy the Liouvillian is
held fixed while the optical field is treated perturbatively. In
response geometry, the Liouvillian itself becomes a
parameter-dependent object. The perturbation is no longer only
the electric field, but motion through the manifold of dynamical
generators. This shift converts nonlinear spectroscopy from a
calculation of response amplitudes into a probe of Liouville-space
transport.
The following sections develop this idea explicitly for open
quantum systems. The response tensor constructed from
stationary-state expectation values yields the intrinsic metric and
curvature fields. The Liouvillian formulation then supplies the
connection and curvature governing pathway transport. Finally,
nonlinear spectroscopy observes a projection of this transport
through a chosen measurement channel, giving rise to the
observable bundle and observational holonomy.


\subsection{Quantum Regression and Liouville-Space Response}
\label{subsec:qrt_response}
The Green--Kubo formulation expresses linear response in terms
of equilibrium two-time correlation functions. For open quantum
systems, the corresponding object is not a Hamiltonian
propagator but the Liouvillian propagator generated by the
master equation. This connection is supplied by the quantum
regression theorem.
Consider a Markovian open quantum system whose reduced
density operator evolves according to
\begin{align}
\frac{d\rho}{dt}
=
\mathcal L\rho .
\end{align}
If the same dynamical semigroup governs both the evolution of
the density operator and the relaxation of conditional
correlations, then the two-time correlation function may be
written as
\begin{align}
\langle A(t)B(0)\rangle_{\rm ss}
=
{\rm Tr}
\left[
A e^{\mathcal L t}
\left(
B\rho_{\rm ss}
\right)
\right].
\label{eq:QRT_two_time}
\end{align}
Equation~\eqref{eq:QRT_two_time} is the quantum regression
theorem. It states that the same Liouvillian propagator that
evolves the density matrix also propagates the operator-valued
initial condition generated by the first measurement or field
interaction.

This result is central for the geometric formulation developed
here. It means that experimentally accessible correlation
functions inherit their parameter dependence from the
Liouvillian itself. If
\begin{align}
\mathcal L=\mathcal L(\lambda),
\end{align}
then the response functions measured in spectroscopy are
functions on the same Liouvillian control manifold:
\begin{align}
\langle A(t)B(0)\rangle_{\lambda}
=
{\rm Tr}
\left[
A e^{\mathcal L(\lambda)t}
\left(
B\rho_{\rm ss}(\lambda)
\right)
\right].
\end{align}
Thus, the geometry of the Liouvillian is directly imprinted onto
the geometry of measured response functions.

The same structure extends to multi-time correlations. For
three operators,
\begin{align}
\langle A(t_2+t_1)B(t_1)C(0)\rangle_{\rm ss}
=
{\rm Tr}
\left[
A e^{\mathcal L t_2}
B e^{\mathcal L t_1}
\left(
C\rho_{\rm ss}
\right)
\right],
\end{align}
with the ordering of operators and propagators determined by
the experimental sequence. Higher-order nonlinear response
functions are built from the same ingredients: field-interaction
superoperators separated by Liouvillian propagators. For
example, a third-order optical response can be written
schematically as
\begin{align}
R^{(3)}(t_3,t_2,t_1)
=
{\rm Tr}
\left[
\mu_-
e^{\mathcal L t_3}
\mu_\times
e^{\mathcal L t_2}
\mu_\times
e^{\mathcal L t_1}
\mu_\times
\rho_{\rm ss}
\right],
\label{eq:R3_liouville_qrt}
\end{align}
where \(\mu_\times\rho=[\mu,\rho]\) denotes the
light--matter interaction superoperator and \(\mu_-\) represents
the detection projection.
Equation~\eqref{eq:R3_liouville_qrt} shows why the Liouvillian
rather than the Hamiltonian is the natural geometric object for
spectroscopic response. The nonlinear signal is not determined
only by eigenenergies or transition dipoles. It is determined by
successive propagation through Liouville space. Consequently,
any deformation of the Liouvillian deforms the entire
multi-time response function.

The quantum regression theorem therefore provides the
mathematical bridge between conventional response theory and
observational geometry. Kubo theory identifies response
functions with correlation functions. The regression theorem
identifies those correlation functions with Liouvillian
propagation. Response geometry then studies how these
Liouvillian propagators, and hence the measured response
functions, change as the generator moves through a manifold of
control parameters.

Having established that spectroscopy measures functions built
from \(e^{\mathcal L t}\), the next step is to ask how those
functions change when \(\mathcal L\) is deformed. This is
precisely the role of the Duhamel expansion.


%-----
